% Port aus SwissTeX 1.3.1 (Produktionsfork, test-display.tex), Umschlag auf
% die v2.0-Slot-Form gehoben: die alte Positionsform
% \swisscover{K}{T}{U}{F} plus \renewcommand{\swisscoverglyph}{S} wird zu
% \swisscover*{kicker=..., ..., glyph=S} -- genau der Uebergang, fuer den
% der Abkuendigungsmechanismus (swisstex.cls Abschnitt 10a) gebaut ist. Die
% Vorlage uebt bewusst die MODERNE Form aus, siehe
% tests/test_regression_docs.py::test_display_build_is_warning_free.
\documentclass[tint=true]{swisstex}
\setrunningtitle{SwissTeX \textbar{} Anzeigengrade und Farbfelder}
\begin{document}

\swisscover*{kicker=Musterblatt, title=SwissTeX,
  subtitle={Anzeigengrade, Farbfelder, Gliederung},
  foot={TeX Gyre Heros \textbar{} Raster 13,5\,pt \textbar{} Satzspiegel 105\,mm},
  glyph=S}

\section{Anzeigengrade}

\subsection{Vom Raster abgeleitet}

Jeder Schaugrad hat einen Durchschuss von ganzen Rasterzeilen. Der
Schriftgrad folgt daraus, grosse Grade laufen enger.

\begin{gridblock}
{\swissdisplay{4}\bfseries Raster}\par
{\swissdisplay{3}Grotesk}\par
{\swissdisplay{2}Satzspiegel}\par
{\swissdisplay{1}Basislinienraster}
\end{gridblock}

Nach den Schaugraden setzt der Grundtext wieder auf dem Raster auf.

\section{Farbfeld}

\subsection{Band über das volle Mass}

\begin{swissband}
{\swissdisplay{2}\bfseries Kunsthalle}\par\vspace{6pt}
Ein Farbfeld läuft über Marginalspalte und Satzspiegel. Der Satz steht in
Weiss, die Höhe bleibt am Raster.
\end{swissband}

Nach dem Band setzt der Grundtext wieder auf dem Raster auf, ohne Versatz.

\subsection{Zweiter Absatz}

Damit lässt sich prüfen, ob die Rasterbindung über mehrere Blöcke hinweg
hält.

\end{document}
